\documentclass[12pt]{article}
\usepackage[utf8]{inputenc}
\usepackage[T1]{fontenc}
\usepackage{hyperref}
\title{Sori-Harengan Language Resources for K-12 Curriculum Development}
\author{Compiled by Claude Sonnet 5.5 (claude-sonnet-5-5) for Andr\'as Kornai}
\date{}
\begin{document}
\maketitle

\section{Introduction}
Sori-Harengan (ISO 639-3 sbh, Glottocode sori1242) is a West Manus Austronesian language of Manus Province, Papua New Guinea. It is spoken on the Sori and Harengan islands and the northwest coast of Manus Island. Sources give about 570 speakers (1977) and about 700 (Blust, 2021). Joshua Project reports about 2,200 for the people group. The word order is SVO. Searches found no school materials, dictionary or grammar.

\section{Primary Educational Resources}
No primers, textbooks or school programmes were found in search results.

\section{Linguistic and Reference Materials}
\begin{itemize}
\item Wikipedia overview, which notes that Blust's Sori-Harengan notes are archived with Kaipuleohone: \url{https://en.wikipedia.org/wiki/Sori-Harengan_language}
\item Numeral system data (Blust, from the Sori dialect), Max Planck EVA: \url{https://lingweb.eva.mpg.de/channumerals/Sori-Harengan.htm}
\item ASJP wordlist (compiled from Z'graggen 1975): \url{https://asjp.clld.org/languages/SORI_HARENGAN}
\end{itemize}

\section{Religious and Community Resources}
\begin{itemize}
\item Joshua Project profile (Scripture translation listed as started): \url{https://joshuaproject.net/people_groups/15002}
\item Scripture Earth entry for sbh: \url{https://www.scriptureearth.org/00eng.php?iso=sbh}
\end{itemize}
I did not verify the contents of the Scripture Earth page.

\section{Digital and Audio Resources}
\begin{itemize}
\item Global Recordings Network, Sori (varieties Sori and Harengan; Tok Pisin is also understood): \url{https://globalrecordings.net/en/language/sbh}
\end{itemize}

\section{Curriculum Development Context}
The language is likely spoken by the whole community, but it is not known to be taught in schools. No orthography or classroom materials were found. A curriculum would first need an agreed orthography and a basic lexicon, using Tok Pisin and English as bridging languages.
\end{document}